% Source: https://artofproblemsolving.com/community/c3h2320771_1000th_post_problems
% Stable identity: aops-topic-2320771-q05
Let the angle be $x^\circ$ with $0<x<90$. Its supplement is $180-x$ and its complement is $90-x$, so
\[180-x=10(90-x)\iff9x=720\iff x=\boxed{80}.\]

\begin{asy}
usepackage("amsmath");
size(4cm);
defaultpen(fontsize(10pt));
pair O=(0,0);
draw(O--2.2*dir(0),linewidth(0.9)); draw(O--2.2*dir(80),linewidth(0.9));
draw(O--2.2*dir(90),dashed); draw(O--2.2*dir(180),dashed);
draw(arc(O,0.6,0,80),blue); label("$80^\circ$",0.85*dir(30),blue);
draw(arc(O,1.4,80,90),red); label("$10^\circ$",1.45*dir(80),E,red);
draw(arc(O,1.0,80,180),deepgreen); label("$100^\circ$",1.2*dir(135),deepgreen);
dot(O); label("$O$",O,S);
\end{asy}

Check: the complement of $80^\circ$ is $10^\circ$ and its supplement is $100^\circ=10\cdot10^\circ$.
